% =============================================================================
% Referee patch: §7 Actual-zeta transverse suppression
% Purpose: make the Section 5 -> Section 7 scalar-functional handoff type-safe.
%
% This patch does not prove actual-zeta suppression.  It only fixes what the
% Section 7 hypotheses mean after the new scalar toy anomaly functional
% A_toy(X)=X_{12}+X_{21} is adopted.
% =============================================================================

% =============================================================================
\section{Actual-zeta transverse suppression}
\label{sec:actual-zeta-transverse-suppression}
\statuspaper{LIVE}\quad
\statusmig{not-started}\quad
\statusreferee{in-progress}\quad
\statussympy{not-started}\quad
\statussim{not-started}\quad
\statuslean{not-started}
\par\medskip

The actual zeta phase \(\Ph\) defines two whitened blocks --- the
real-axis two-center block \(\widehat\Om_\z(s;m)\) and a four-center
analogue \(\widehat\Om_\z^{(4)}\) for the symmetry quartet.  The scalar
functional used on the actual-zeta side must be the same scalar channel
selected on the toy side, namely the off-diagonal trace
\[
        \mathcal A_\toy(X)=X_{12}+X_{21}.
\]
If the transverse projection of Section~\ref{sec:transverse-projection}
has already been constructed, this sentence is read as follows: there is
a scalar functional \(\ell_2\) on the two-point transverse quotient such
that
\begin{equation}
\label{eq:A2-transverse-representative}
        \mathcal A_2(X):=\ell_2(\Pi_\trans X)
\end{equation}
agrees with \(X_{12}+X_{21}\) on the chosen normalized representatives.
Before Section~\ref{sec:transverse-projection} is closed, \(\mathcal A_2\)
should simply be read on representatives as
\begin{equation}
\label{eq:A2-representative-functional}
        \mathcal A_2(X)=X_{12}+X_{21}.
\end{equation}
Thus \(\Pi_\trans\) acts on the matrix or its quotient class before any
scalar evaluation; it is not applied after the scalar has already been
taken.

The suppression hypotheses below are analytic hypotheses about the
actual-zeta blocks in this fixed scalar channel.  They do not use the toy
phase, except through the choice of scalar channel, the retained compact
separation set \(D\) from Theorem~\ref{thm:toy-anomaly-visibility}, and
the toy exponent \(A_\toy=0\).  In particular, changing the toy phase in
Section~\ref{sec:toy-anomaly-and-transverse-obstruction} does not change
the actual-zeta block \(\widehat\Om_\z\); it changes only which scalar
channel and which toy lower-bound constant are compared against the
actual-zeta upper bound.

\subsection{Two-point suppression}
\label{subsec:two-point-suppression}

\begin{hypothesis}[Two-point actual-zeta scalar suppression]
\label{hyp:zeta-suppression-two-point}
Let \(D\subset(0,\infty)\) be the compact separation interval retained
in Theorem~\ref{thm:toy-anomaly-visibility}.  For the two-center
actual-zeta packet associated with \(\rho\) and a symmetry partner, after
the gauge-fixed normalization of Definition~\ref{def:local-phase-chart}
and after the transverse representative convention of
\eqref{eq:A2-transverse-representative}--\eqref{eq:A2-representative-functional},
\begin{equation}
\label{eq:zeta-two-point-scalar-suppression}
        \bigl|\mathcal A_2(\widehat\Om_\z(s;m))\bigr|
        \le Q^{-B_{\z,2}}
\end{equation}
on the retained subfamily with dimensionless separation \(d\in D\), for
an absolute exponent \(B_{\z,2}>A_\toy=0\).  Here \(\mathcal A_2\) is
the actual-zeta realization of the same scalar channel
\(X_{12}+X_{21}\) used in
Definition~\ref{def:toy-anomaly-functional}.
\end{hypothesis}

\begin{wip}\label{rem:wip-zeta-suppression-two-point}
The two-point case of the structural gap.  Consumed by
Theorem~\ref{thm:two-point-rigidity}.  Admissible proof routes:
\textup{(i)} a standalone route from the local jet-Gram geometry alone;
\textup{(ii)} an energy-supply route consuming the source bound
of Section~\ref{sec:source-energy-supply} via
Lemma~\ref{lem:source-energy-to-suppression}.  Existing v1 scalar
estimates involving
\(\Phi_K(X)=X_{12}+X_{21}\), such as estimates for
\(H_m(s)=\Phi_K(\widehat\Omega_\zeta^{\mathrm{corr}}(s;m))\), are
candidates for import only after their normalization, correction package,
and representative convention are matched to
\eqref{eq:A2-transverse-representative}.  V1 results using the toy
matrix \(M(d)\), \(A_{\val}\), or \(\Psi_{x,d}\) are not imported by
this hypothesis.
\end{wip}

\subsection{Four-point suppression}
\label{subsec:four-point-suppression}

\begin{hypothesis}[Four-point actual-zeta scalar suppression]
\label{hyp:zeta-suppression-four-point}
For the four-center packet associated with the symmetry quartet
\(\{\rho,\bar\rho,1-\rho,1-\bar\rho\}\), let
\(\mathcal A_4\) denote the four-point scalar channel chosen as the
four-point analogue of \(\mathcal A_2\), after the four-point quotient
representative has been fixed.  Then
\begin{equation}
\label{eq:zeta-four-point-scalar-suppression}
        \bigl|\mathcal A_4(\widehat\Om_\z^{(4)})\bigr|
        \le Q^{-B_{\z,4}}
\end{equation}
on the retained four-point subfamily, for an absolute exponent
\(B_{\z,4}>A_\toy^{(4)}\).
\end{hypothesis}

\begin{wip}\label{rem:wip-zeta-suppression-four-point}
The four-point case of the structural gap.  Consumed by
Theorem~\ref{thm:four-point-rigidity}.  Same admissible proof routes
as Hypothesis~\ref{hyp:zeta-suppression-two-point}.  The scalar channel
\(\mathcal A_4\) must be fixed before the proof of four-point
suppression and must not be chosen after seeing the actual-zeta error
terms.  Any v1 four-point import must be checked at the scalar-channel
level and may not silently use the archived two-point toy matrix
calibration.
\end{wip}

\begin{remark}[Safe scalar import from v1]
\label{rem:safe-scalar-import-from-v1}
The new toy phase of Section~\ref{sec:toy-anomaly-and-transverse-obstruction}
changes the toy lower-bound calculation, but it does not change the
actual-zeta phase or the scalar functional \(X_{12}+X_{21}\).  Therefore
v1 estimates that are purely estimates of the scalar
\(\Phi_K(\widehat\Omega_\zeta^{\mathrm{corr}})\), with
\(\Phi_K(X)=X_{12}+X_{21}\), are structurally compatible with the
present rebuild.  By contrast, any v1 argument that normalizes by the
archived toy matrix \(M(d)\), pairs with \(A_{\val}\), or invokes the
calibration map \(\Psi_{x,d}\), is not compatible without a fresh
rederivation for the new toy.
\end{remark}

\begin{remark}[Standalone and energy-supply routes]
\label{rem:standalone-and-energy-supply-routes}
Hypotheses~\ref{hyp:zeta-suppression-two-point}
and~\ref{hyp:zeta-suppression-four-point} are stated independently of
the technology of their proofs.  A purely local proof would use only
Sections~\ref{sec:jet-limit-local-blocks}--\ref{sec:transverse-projection};
an energy-supply proof may additionally use
Section~\ref{sec:source-energy-supply}.  In either route, the output must
be an upper bound in the fixed scalar channel appearing in
\eqref{eq:zeta-two-point-scalar-suppression} or
\eqref{eq:zeta-four-point-scalar-suppression}; source-energy estimates in
other channels are not sufficient unless transferred into this channel.
\end{remark}
